
\chapter{深度与\trCM{}性}
\label{ch:depth}

我们现在引入\trDepth{}的观念。深度是一个代数性而非几何性的概念，因此我们把精力放在培养「它意味着什么」的几何直觉上。最重要的是（诸）伴点（的闭包）之维数对深度的几何制约（\mbox{定理~26.1.3}）。理解深度的核心工具是 Koszul 复形（\S24.4.1 中略作讨论），但我们将绕开这条路径，因为所需之物都能直接证明。
当局部环的深度等于它的维数时，该环便称作 \trCM{}的，而这是概形得以「好」的一种重要方式。例如，正则局部环是\trCM{}的（\S26.2.5），光滑簇中的正则嵌入亦然（\mbox{命题~26.2.6}）。\trCM{}性将是第 29 章\trSerreDuality{}证明的关键，其桥梁是\mbox{\trMiracleFlat{}定理~26.2.11}。
深度的另一项应用是 Serre 的 $R_1+S_2$ 正规性判别法（\mbox{定理~26.3.2}）：我们将用它证明正则概形是正规的（\S26.3.5），而无须先证出「正则局部环都是因子分解的」（\mbox{Auslander--Buchsbaum 定理，事实~13.8.5}）；还将用它证明：光滑概形中的正则嵌入若在余维数 $1$ 处正则，则是正规的（\S26.3.3）。


\section{深度}

\textbf{26.1.1.}\ \emph{非正式导言。}\quad 给定\trScheme{} $X$ 上一点 $\mathrm{p}$，有时你可以局部地「切」过这个点——也就是说，存在一个经过 $\mathrm{p}$ 的\trEffectiveCartierDivisor（比如说在 $\mathrm{p}$ 的某个\trNeighborhood 中）。你能否这样做，当且仅当这个点\emph{不是} $X$ 的\trAssociatedPoint。有时你还能再切一刀：只要切出的第一片仍不以 $\mathrm{p}$ 为\trAssociatedPoint。（$X$ 在点 $\mathrm{p}$ 处）的\trDepth{}说的正是你能切多深——即在这个点变成\trAssociatedPoint{}而把你挡住之前，你总共能切出多少刀。显然切数不可能超过 $\dim X$：因为每切一刀便失去一维，而切满 $\dim X$ 刀之后，（在集合论意义下）就只剩下点 $\mathrm{p}$ 了（此时它便是一个\trAssociatedPoint）。于是\trDepth{}以\trDimension{}为上界。在求\trDepth{}时，你或许以为必须小心地切，才能让切出的刀数最多。不必担心：随手乱切，最后也总会停在同一步上。

\textbf{26.1.2.}\ \emph{\trhDefinition。}\quad 设 $(A,\mathfrak{m})$ 是一个\trNoetherian\trLocalRing，$M$ 是一个\trFinitelyGenerated{}的 $A$-\trModule。$M$ 的\textbf{\trDepth}（记作 $\operatorname{depth}M$）是取值于 $\mathfrak{m}$ 中的最长 $M$-\trRegularSequence{}的\trLength。（\trRegularSequence{}的理论曾在 \S9.5.3 中讨论。）

（更一般地，若 $R$ 是\trNoetherian{}环，$I\subset R$ 是一个\trIdeal，而 $M$ 是一个\trFinitelyGenerated{}的 $R$-\trModule，则 $M$ 的\textbf{$I$-\trDepth}（记作 $\operatorname{depth}_I M$）是取值于 $I$ 中的最长 $M$-\trRegularSequence{}的\trLength。我们用不到这个概念。）

\textbf{26.1.A.（重要练习）}\quad 设 $M$ 是\trNoetherian\trLocalRing{} $(A,\mathfrak{m})$ 上一个非零\trFinitelyGenerated{}的\trModule。证明：$\operatorname{depth}M=0$，当且仅当 $\mathfrak{m}$ 的每个元素都是 $M$ 的\trZerodivisor，当且仅当 $\mathfrak{m}$ 是 $M$ 的一个\trAssociatedPrimeIdeal。

\textbf{26.1.B.（练习）}\quad 设 $M$ 是\trNoetherian\trLocalRing{} $A$ 上一个\trFinitelyGenerated{}的\trModule。证明：$\operatorname{depth}M\le\dim\Supp M$。特别地，
\[
\tag{26.1.2.1}\label{eq:depth-le-dim}
\operatorname{depth}A\le\dim A.
\]
（我们将在\mbox{定理~26.1.7} 中改进这个结果。）

到了这一步，要确定一个 $A$-\trModule{} $M$ 的\trDepth{}仍然很难。你可以试着依次选取 $x_1,x_2,\dots$ 来构造一条 $M$-\trRegularSequence，可你怎么知道自己做出的选取恰好给出了\emph{最长}的那条？令人高兴的答案是：你不可能出错；这正是下一条结果的内容。（随后我们将在 \S26.1.5 中描述实际操作中如何求出一个\trModule{} $M$ 的\trDepth。）

\textbf{26.1.3. \trhTheorem。}\ ——\ \emph{设 $M$ 是\trNoetherian\trLocalRing{} $(A,\mathfrak{m})$ 上一个\trFinitelyGenerated{}的\trModule。那么含于 $\mathfrak{m}$ 中的所有极大 $M$-\trRegularSequence（即 $\mathfrak{m}$ 中那些不能延长为 $x_1,\dots,x_{n+1}$ 的 $M$-\trRegularSequence{} $x_1,\dots,x_n$）都有相同的\trLength。于是 $M$ 的\trDepth{}就是任何极大 $M$-\trRegularSequence{}的\trLength。}

我们将给出刻画「\trRegularSequence{}何时极大」的一个\trCohomology{}判别法（\mbox{定理~26.1.4}），以此证明\mbox{定理~26.1.3}。（然后你将在\mbox{练习~26.1.D} 中证明\mbox{定理~26.1.3}。）我们还将用这个判别法，在\mbox{定理~26.1.7} 中给出对\trDepth{}的一个更好的界。\mbox{定理~26.1.3} 是本章关键的技术性结果。一条重要的寓意是：应当把\trDepth{}理解为一种「\trCohomology{}」性质。

\textbf{26.1.4. \trhTheorem{}（\trRegularSequence{}存在性的\trCohomology{}判别法）。}\ ——\ \emph{设 $(A,\mathfrak{m})$ 是一个\trNoetherian\trLocalRing，$M$ 是一个非零\trFinitelyGenerated{}的 $A$-\trModule。则以下诸条等价。}

(i) \emph{对每个满足 $\Supp N=\{[\mathfrak{m}]\}$ 的\trFinitelyGenerated{} $A$-\trModule{} $N$（即 $N$ 具有\trFiniteLength，\mbox{练习~6.6.X}(a)），当 $i<n$ 时都有 $\mathrm{Ext}^i_A(N,M)=0$。}

(ii) \emph{当 $i<n$ 时都有 $\mathrm{Ext}^i_A(A/\mathfrak{m},M)=0$。}

(iii) \emph{在 $\mathfrak{m}$ 中存在一条长度为 $n$ 的 $M$-\trRegularSequence。}

这个结果可以用多种方式推广；例如参看 [Mat1, Thm.~28]。

\emph{证明。}\quad 显然 (i) 蕴含 (ii)。

\textbf{26.1.C.（练习）}\quad 证明 (ii) 蕴含 (i)。提示：对 $N$ 的\trLength{}作归纳，并用 $\mathrm{Ext}^i_A(\cdot,M)$ 的\trLongExactSequence。

\emph{(iii) 蕴含 (ii) 的证明。}\quad $n=0$ 的情形是空洞的。我们对一切 $n$ 归纳地证明该结果。设 (iii) 在 $n\ge 1$ 时满足，并假设已知 (ii) 对「所有更小的 $n$」成立。取一条长度为 $n$ 的\trRegularSequence{} $x_1,\dots,x_n$。于是 $x_1$ 是 $M$ 上的一个\trNonZerodivisor，故我们有\trExactSequence{}
\[
\tag{26.1.4.1}\label{eq:depth-regular-ses}
0 \longrightarrow M \overset{\times x_1}{\longrightarrow} M \longrightarrow M/x_1 M \longrightarrow 0.
\]
于是 $M/x_1 M$ 有长度为 $n-1$ 的\trRegularSequence{} $x_2,\dots,x_{n-1}$，故由归纳假设，当 $i<n-1$ 时 $\mathrm{Ext}^i_A(A/\mathfrak{m},M/x_1 M)=0$。对 \eqref{eq:depth-regular-ses} 取 $\mathrm{Ext}^i_A(A/\mathfrak{m},\cdot)$ 的 Ext \trLongExactSequence，我们发现
\[
\mathrm{Ext}^i_A(A/\mathfrak{m},M) \overset{\times x_1}{\longrightarrow} \mathrm{Ext}^i_A(A/\mathfrak{m},M)
\]
当 $i<n$ 时是\trInjectiveMap。现在，$\mathrm{Ext}^i_A(A/\mathfrak{m},M)$ 可以通过取 $M$ 的一个\trInjective{}分解并施加 $\mathrm{Hom}_A(A/\mathfrak{m},\cdot)$ 来计算。于是，由于 $x_1$ 落在 $\mathfrak{m}$ 中（从而零化 $A/\mathfrak{m}$），乘以 $x_1$ 的映射是零映射。故 (ii) 对 $n$ 同样成立。

\emph{(ii) 蕴含 (iii) 的证明。}\quad $n=0$ 的情形是空洞的。

接下来我们处理 $n=1$ 的情形，办法是证明其逆否命题。假设 $\mathfrak{m}$ 中没有 $M$ 上的\trNonZerodivisor，于是由\mbox{练习~26.1.A}，$\mathfrak{m}$ 是 $M$ 的一个\trAssociatedPrimeIdeal。这样（\S6.6.7）我们有一个\trInjectiveMap{} $A/\mathfrak{m}\hookrightarrow M$，从而 $\mathrm{Hom}_A(A/\mathfrak{m},M)\neq 0$，正合所需。

现在我们对一切 $n>1$ 归纳地证明该结果。设 (ii) 在 $n\ge 2$ 时满足，并假设已知 (iii) 对「所有更小的 $n$」成立。那么由 $n=1$ 的情形，存在 $M$ 上的一个\trNonZerodivisor{} $x_1$，于是我们有\trShortExactSequence{} \eqref{eq:depth-regular-ses}。对 \eqref{eq:depth-regular-ses} 施加 $\mathrm{Ext}^i_A(A/\mathfrak{m},\cdot)$ 所得 Ext \trLongExactSequence{}的一个片段是
\[
\mathrm{Ext}^i_A(A/\mathfrak{m},M) \longrightarrow \mathrm{Ext}^i_A(A/\mathfrak{m},M/x_1 M) \longrightarrow \mathrm{Ext}^{i+1}_A(A/\mathfrak{m},M).
\]
由假设，当 $i<n-1$ 时 $\mathrm{Ext}^i_A(A/\mathfrak{m},M)$ 与 $\mathrm{Ext}^{i+1}_A(A/\mathfrak{m},M)$ 均为 $0$，故当 $i<n-1$ 时 $\mathrm{Ext}^i_A(A/\mathfrak{m},M/x_1 M)=0$，于是由归纳假设，我们有一条在 $\mathfrak{m}$ 中长度为 $n-1$ 的 $(M/x_1 M)$-\trRegularSequence{} $x_2,\dots,x_n$。把 $x_1$ 添加到这条序列的最前端，即得所证。\hfill$\square$

\textbf{26.1.D.（练习）}\quad 证明\mbox{定理~26.1.3}。提示：由\mbox{定理~26.1.4}（尤其是 (ii) 与 (iii) 的等价性），你已经掌控了 $\mathfrak{m}$ 中的 $M$-\trRegularSequence{}至多能有多长。请使用这个判别法，以及\mbox{定理~26.1.4} 的证明中所用的那条\trLongExactSequence，证明 $\mathfrak{m}$ 中的任何 $M$-\trRegularSequence{}都能延长到这个长度。

\textbf{26.1.E.（练习）}\quad 设 $M$ 是\trNoetherian\trLocalRing{} $(A,\mathfrak{m})$ 上一个\trFinitelyGenerated{}的\trModule。若 $x$ 是 $\mathfrak{m}$ 中的一个\trNonZerodivisor，证明：$\operatorname{depth}(M/xM)=\operatorname{depth}M-1$。提示：\mbox{定理~26.1.3}。

\textbf{26.1.5.}\ \emph{求\trModule{}的\trDepth。}\quad 我们现在可以按如下方式，通过逐个找出\trNonZerodivisor{}来计算 $M$ 的\trDepth。在 $\mathfrak{m}$ 中是否存在 $M$ 上的\trNonZerodivisor{} $x$？

(a) 若不存在，则 $\operatorname{depth}M=0$（\mbox{练习~26.1.A}）。

(b) 若存在，则任取这样一个 $x$，并（用上一道练习）对 $M/xM$ 重复这一过程。

由\mbox{练习~26.1.B}，这一过程必定终止。

\textbf{26.1.F.（重要练习）}\quad 设 $(A,\mathfrak{m})$ 是维数为 $d$ 的\trRegular\trLocalRing。证明：$\operatorname{depth}A=d$。提示：\mbox{练习~13.2.B}。

\textbf{26.1.G.（练习（参看\mbox{练习~24.3.I}））}\quad 设 $X=\Spec R$，其中
\[
R = k[w,x,y,z]/(wy,wz,xy,xz),
\]
即 $\mathbb{A}^4_k$ 中相交于原点的两个坐标二维平面的并。证明：$X$ 在原点处的\trLocalRing{}的\trDepth{}为 $1$。提示：证明 $w-y$ 不是\trZerodivisor，而 $R/(w-y)$ 在原点处有一个\trEmbeddedPoint。

\textbf{26.1.6. \trDepth{}以\trAssociatedPrimeIdeal{}的\trDimension{}为上界。}

\mbox{定理~26.1.4} 可用来对「以维数限制\trDepth」的界 \eqref{eq:depth-le-dim} 给出一个重要的改进：

\textbf{26.1.7. \trhTheorem。}\ ——\ \emph{\trModule{} $M$ 的\trDepth{}至多是诸 $\dim A/\mathfrak{p}$ 中的最小者，其中 $\mathfrak{p}$ 取遍 $M$ 的诸\trAssociatedPrimeIdeal。}

（\mbox{练习~26.1.G} 中两张平面交于一点的例子表明，这个界并不锐利。）\mbox{定理~26.1.7} 证明中的关键一步，是下述属于 Ischebeck 的结果。

\textbf{26.1.8. \trhLemma。}\ ——\ \emph{设 $(A,\mathfrak{m})$ 是一个\trNoetherian\trLocalRing，$M$ 与 $N$ 是非零\trFinitelyGenerated{}的 $A$-\trModule。那么当 $i<\operatorname{depth}M-\dim\Supp N$ 时 $\mathrm{Ext}^i_A(N,M)=0$。}

\emph{证明。}\quad 考虑如下诸论断。

(general$_\tau$) \emph{\mbox{\trhLemma~26.1.8} 在 $\dim\Supp N\le\tau$ 时成立。}

(prime$_\tau$) \emph{\mbox{\trhLemma~26.1.8} 在 $\dim\Supp N\le\tau$ 且 $N=A/\mathfrak{p}$（$\mathfrak{p}$ 为某个\trPrimeIdeal）时成立。}

注意（general$_0$）（从而（prime$_0$））为真：因为此时 $\Supp N=\{[\mathfrak{m}]\}$，而该结果由\mbox{定理~26.1.4}（由「(iii) 蕴含 (i)」）推出。

\textbf{26.1.H.（练习）}\quad 证明（prime$_\tau$）蕴含（general$_\tau$）。提示：对\trShortExactSequence{} $0\to M'\to M\to M''\to 0$，有 $\Supp M=\Supp M'\cup\Supp M''$（\mbox{练习~6.6.B}）。

我们最后证明：当 $\tau\ge 1$ 时，（general$_{\tau-1}$）蕴含（prime$_\tau$），以此结束证明。取定一个满足 $\dim A/\mathfrak{p}=\tau$ 的\trPrimeIdeal{} $\mathfrak{p}\subset A$。由于 $\dim A/\mathfrak{p}>0$，我们可以选取 $x\in\mathfrak{m}\setminus\mathfrak{p}$。考虑\trExactSequence{}
\[
\tag{26.1.8.1}\label{eq:depth-quotient-ses}
0 \longrightarrow A/\mathfrak{p} \overset{\times x}{\longrightarrow} A/\mathfrak{p} \longrightarrow A/(\mathfrak{p}+(x)) \longrightarrow 0,
\]
注意 $\dim\Supp A/(\mathfrak{p}+(x))\le\tau-1$（你明白这是为什么吗？）。然后，对 \eqref{eq:depth-quotient-ses} 施加 $\mathrm{Hom}_A(\cdot,M)$ 得到的 Ext \trLongExactSequence，连同 $\mathrm{Ext}^i_A(A/(\mathfrak{p}+(x)),M)$\footnote{原书此处作「$\mathrm{Ext}^i_A(A/(\mathfrak{p}+(x))),M)$」，其中「$A/(\mathfrak{p}+(x))$」之后多印了一枚右括号，当为「$\mathrm{Ext}^i_A(A/(\mathfrak{p}+(x)),M)$」之误，径改。——译者注}在 $i<\operatorname{depth}M-\tau+1$ 时的消失（由（general$_{\tau-1}$）），便蕴含
\[
\mathrm{Ext}^i_A(A/\mathfrak{p},M) \overset{\times x}{\longrightarrow} \mathrm{Ext}^i_A(A/\mathfrak{p},M)
\]
在 $i<\operatorname{depth}M-\tau$ 时是\trIsomorphism。但 $\mathrm{Ext}^i_A(A/\mathfrak{p},M)$ 是一个\trFinitelyGenerated{}的 $A$-\trModule（\mbox{练习~23.2.F}），故由\mbox{\trNakayama{}引理~8.2.8}，当 $i<\operatorname{depth}M-\tau$ 时 $\mathrm{Ext}^i_A(A/\mathfrak{p},M)=0$。\hfill$\square$

\textbf{26.1.I.（容易的练习）}\quad 证明\mbox{定理~26.1.7}。提示：$\mathfrak{p}\in\operatorname{Ass}(M)$ 时 $\mathrm{Hom}_A(A/\mathfrak{p},M)\ne 0$（\S6.6.7）。



\section{\trCM{}环与概形}

\textbf{26.2.1.}\ \emph{\trhDefinition。}\quad \trNoether{}局部环 $(A,\mathfrak{m})$ 称为\textbf{\trCM}的，若 $\operatorname{depth}A=\dim A$，即在 \eqref{eq:depth-le-dim} 中等号成立。（类似地也可以定义\trCM{}模，但我们用不到这个概念。）一个\trLocally\trNoether{}\trScheme{}称为\textbf{\trCM}的，若它的所有\trLocalRing{}都是\trCM{}的。我们不会使用这种说法，不过：环 $A$ 称为\textbf{\trCM{}环}，若它在所有素理想 $\mathfrak{p}$ 处的局部化 $A_{\mathfrak{p}}$ 都是\trCM{}的——等价地，若 $\Spec A$ 是\trCM{}的。（下面的\mbox{事实~26.2.14} 说：局部环的\trCM{}性在局部化下保持。\footnote{原书此处作「Cohen-Maculayness」，脱一字母「a」，当为「Cohen-Macaulayness」之误，径改。——译者注}这便保证了此定义与\trCM{}局部环的定义之间不存在矛盾。）

\textbf{26.2.A.（练习）}\quad 证明：每个维数为 $0$ 的\trLocally\trNoether{}\trScheme{}都是\trCM{}的。再证明：维数为 $1$ 的\trLocally\trNoether{}\trScheme{}是\trCM{}的，当且仅当它没有\trEmbeddedPoint。

\textbf{26.2.2.}\ \emph{（反）例。}\quad 设 $X$ 是\mbox{练习~26.1.G} 的那个例子——交于一点的两个平面。由\mbox{练习~26.1.G}，$X$ 不是\trCM{}的。

\textbf{26.2.B.（练习）}\quad 设 $A$ 是\trCM{}的\trNoether{}局部环。用\mbox{定理~26.1.7} 证明：$\Spec A$ 是\trPureDimensional{}的，并且没有\trEmbeddedPoint。（\trNoether{}局部环即使有\trPureDimension{}且无\trEmbeddedPrime，也未必是\trCM{}的；见\mbox{（反）例~26.2.2}。）

\textbf{26.2.3. \trhTheorem（\trCM{}性的\trSlicing{}判别法）。}\ ——\ \emph{设 $(A,\mathfrak{m})$ 是\trNoether{}局部环，$x\in\mathfrak{m}$ 是一个\trNonZerodivisor。则 $(A,\mathfrak{m})$ 是\trCM{}的，当且仅当 $(A/x,\mathfrak{m})$ 是\trCM{}的。}

把它与正则性及平坦性的\trSlicing{}判别法相比较（分别是\mbox{练习~13.2.C} 与\mbox{定理~24.6.6}）。

\textbf{26.2.4.}\ \emph{\trSlicing{}判别法的几何解释。}\quad 设 $X$ 是\trLocally\trNoether{}\trScheme，$D$ 是一个\trEffectiveCartierDivisor。若 $X$ 是\trCM{}的，则 $D$ 亦然。若 $D$ 是\trCM{}的，则 $X$ 在 $D$ 的各点处都是\trCM{}的。

\textbf{26.2.C.（练习）}\quad 证明\mbox{定理~26.2.3}：用\mbox{定理~26.1.3}，以及「极大\trRegularSequence（在 $\mathfrak{m}$ 中的）都具有相同长度」这一事实。（提示：回想\mbox{练习~12.3.I}：对任意 $f\in A$ 都有 $\dim A/(f)\ge\dim A-1$。）

\textbf{26.2.D.（练习）}\quad 证明：若 $(A,\mathfrak{m})$ 是\trCM{}的，则一组元素 $x_1,\dots,x_r\in\mathfrak{m}$ 构成（$A$ 的）\trRegularSequence，当且仅当 $\dim A/(x_1,\dots,x_r)=\dim A-r$。

\textbf{26.2.5.}\quad 由\mbox{练习~26.1.F}，\trRegular\trLocalRing（按定义是诺特的）都是\trCM{}的。特别地，由于域 $k$ 上的\trSmooth{}概形是正则的（\mbox{定理~13.2.7(b)}，其证明见 \S24.8.9），我们看到光滑 $k$-\trScheme{}都是\trCM{}的。（结合\mbox{练习~26.2.B}，这便解释了光滑\trVariety{}上的\trEffectiveCartierDivisor{}为何没有\trEmbeddedPoint，印证了\mbox{题外话~20.2.1} 中的一处评注。）将此事与\mbox{练习~26.2.D}（或 \S26.2.4）相结合，便得到下述结论。

\textbf{26.2.6. \trhProposition。}\ ——\ \emph{\trSmooth{} $k$-\trScheme{}中的\trRegularEmbedding{}都是\trCM{}的。}

\textbf{26.2.7.}\quad 作为\mbox{命题~26.2.6} 与「\trCM{}概形没有\trEmbeddedPoint」（\mbox{练习~26.2.B}）的推论，我们看到：\trSmooth{} $k$-\trScheme{}（比如 $\mathbb{A}^n_k$ 或 $\mathbb{P}^n_k$）中的\trRegularEmbedding{}都没有\trEmbeddedPoint，这推广了\mbox{练习~6.6.F}（$\mathbb{A}^n_k$ 中超曲面的情形）。若没有\trCM{}性理论，这一点并不显然！

\textbf{26.2.8.}\ \emph{\trCM{}性的另一种定义。}\quad \trSlicing{}判别法（\mbox{定理~26.2.3}）给出了一种富有启发性的\trCM{}性归纳定义，改用\trEffectiveCartierDivisor{}来表述，其精神与 \S26.1.5 计算\trDepth{}的方法如出一辙。暂且设 $(A,\mathfrak{m})$ 是\trNoether{}局部环。
\begin{itemize}
\item[（i）] 若 $\dim A=0$，则 $A$ 是\trCM{}的（\mbox{练习~26.2.A}）。
\item[（iia）] 若 $\dim A>0$，而 $\mathfrak{m}$ 的每个元素都是\trZerodivisor，则 $A$ 不是\trCM{}的（由\mbox{练习~26.1.A}）。
\item[（iib）] 否则，在 $\mathfrak{m}$ 中任取一个\trNonZerodivisor{} $x$。则 $A$ 是\trCM{}的，当且仅当 $A/(x)$（由\mbox{\trKrull{}\trPrincipalIdealTheorem~12.3.3}，其维数必为 $\dim A-1$）是\trCM{}的。
\end{itemize}

下面这个例子，其实早在我们知道任何代数几何之前就可以陈述出来（只是无法证明）。（我们取 $\mathbb{C}$ 而非任意的域来工作，只是为了让这个命题所需的背景知识尽可能少。）

\textbf{26.2.E.（趣味练习：马克斯·诺特（Max Noether）的 AF$+$BG 定理）}\quad 设 $f,g\in\mathbb{C}[x_0,x_1,x_2]$ 是两个\trHomogeneous{}多项式，它们切出 $\mathbb{P}^2_{\mathbb{C}}$ 中的两条曲线，并「横截地」相交，即交于有限多个既约点。设 $h\in\mathbb{C}[x_0,x_1,x_2]$ 是在这些点处取零值的\trHomogeneous{}多项式。证明：$h\in(f,g)$。提示：证明 $\mathbb{A}^3$ 中仿射锥 $V(f)$ 与 $V(g)$ 的交没有\trEmbeddedPoint。（不借助本章发展的理论，这道题做起来相当不平凡！为示此事之精微：你很容易构造三个二次型 $e,f,g\in\mathbb{C}[x_0,x_1,x_2]$，使它们恰好切出 $\mathbb{P}^2_{\mathbb{C}}$ 中的两个点 $[1,0,0]$ 与 $[0,1,0]$，然而直线 $z=0$ 却因次数原因不在理想 $(e,f,g)$ 之中。）

作为\trCM{}性的又一重动机：下一道练习看似基础，但要把证明做得干净，我们需要\trCM{}性。

\textbf{26.2.F.（练习）}\quad 设 $H_1,\dots,H_n$ 是 $\mathbb{P}^n_{\mathbb{C}}$ 中次数分别为 $d_1,\dots,d_n$ 的 $n$ 张\trHypersurface，并设它们的\trSchemeTheoretic{}交 $Z$ 是有限个点。证明：$Z$ 在诸点处的长度之和（对 $Z$ 的诸点求和）为 $d_1\cdots d_n$。等价地：证明 $Z$ 的\trLength{}为 $d_1\cdots d_n$。不严格地说：这些\trHypersurface{}相交于 $d_1\cdots d_n$ 个点，「按正确方式计数」。提示：\trBezout{}定理，取\mbox{练习~18.6.J} 的形式。

\textbf{26.2.9. \trMiracleFlat。}

我们以一条惊人地简洁而有用的平坦性判别法来收束本节，我们将在\trSerreDuality{}的证明中用到它。其核心内容是下述代数结果。

\textbf{26.2.10. \trhTheorem（\trMiracleFlat{}定理，代数版本）。}\ ——\ \emph{设 $\phi\colon(B,\mathfrak{n})\to(A,\mathfrak{m})$ 是\trNoether{}局部环之间的一个（局部）同态，使得 $A$ 是\trCM{}的，$B$ 是正则的，而 $A/\mathfrak{n}A=A\otimes_B(B/\mathfrak{n})$（对应于该纤维的环）具有纯维数 $\dim A-\dim B$。则 $\phi$ 是平坦的。}

\emph{证明。}\quad 我们对 $\dim B$ 作归纳来证明\mbox{定理~26.2.10}。若 $\dim B=0$，则 $B$ 是一个域（\mbox{练习~13.2.A}），于是结论立得，因为域上的任何东西都平坦（\mbox{练习~24.1.A(a)}）。再设 $\dim B>0$，并假设结论已对一切「维数更小的 $B$」得证。取 $x\in\mathfrak{n}\setminus\mathfrak{n}^2$，于是 $B/x$ 是维数为 $\dim B-1$ 的\trRegular\trLocalRing（\mbox{练习~12.3.I}）。那么
\begin{align*}
\dim A/xA &\le \dim B/(x)+\dim A/\mathfrak{n}A &&\text{（关键练习~12.4.A）}\\
&= \dim B-1+\dim A/\mathfrak{n}A &&\text{（练习~12.3.I）}\\
&= \dim A-1 &&\text{（由定理~26.2.10 的假设）}.
\end{align*}
由（\mbox{练习~12.3.I}）可知 $\dim A/xA\ge\dim A-1$，于是我们有 $\dim A/xA=\dim B/(x)+\dim A/\mathfrak{n}A=\dim A-1$。由\mbox{练习~26.2.D}，$A/xA$ 是\trCM{}环，而 $x$ 在 $A$ 上不是\trNonZerodivisor。归纳假设随之适用于 $(B/(x),\mathfrak{n})\to(A/xA,\mathfrak{m})$，故 $A/xA$ 在 $B/(x)$ 上平坦。再由平坦性的局部\trSlicing{}判别法（在靶上，\mbox{定理~24.6.6}），$B\to A$ 是平坦的，此即所欲证。\hfill$\square$

\textbf{26.2.11. \trhTheorem（\trMiracleFlat{}定理）。}\ ——\ \emph{设 $\pi\colon X\to Y$ 是\trPureDimensional{}的\trFiniteType{} $k$-\trScheme{}之间的一个\trMorphism，其中 $X$ 是\trCM{}的，$Y$ 是正则的，而 $\pi$ 的诸\trFiber{}都具有维数 $\dim X-\dim Y$。则 $\pi$ 是平坦的。}

\textbf{26.2.G.（练习）}\quad 证明\trMiracleFlat\mbox{定理~26.2.11}。（别忘了概形通常还有非闭点！）

\trMiracleFlat\mbox{定理~26.2.11} 中的几何局面，是下面这个漂亮「套装」（package）的一部分。

\textbf{26.2.H.（练习：\trMiracleFlat{}套装）}\quad 设 $\pi\colon X\to Y$ 是\trLocally\trFiniteType{} $k$-\trScheme{}之间的一个映射，其中 $X$ 与 $Y$ 都是\trPureDimensional{}的，且 $Y$ 是正则的。证明：下列三条中若任何两条成立，则第三条也成立：
\begin{itemize}
\item[（i）] $\pi$ 是相对维数 $\dim X-\dim Y$ 的平坦\trMorphism。
\item[（ii）] $X$ 是\trCM{}的。
\item[（iii）] 每个\trFiber{} $X_y$ 都是\trCM{}的，且具有纯维数 $\dim X-\dim Y$。
\end{itemize}
提示：若 $\phi\colon B\to A$ 是平坦的环映射，则 $\phi$ 把\trNonZerodivisor{}送到\trNonZerodivisor（\mbox{观察~24.1.2}）。

\mbox{练习~26.2.H} 的陈述还可以改进，代价是牺牲对称性。在「（ii）与（iii）蕴含（i）」这一方向中，（iii）里的\trCM{}假设并不需要。

\textbf{26.2.12.}\ \emph{\trhExample。}\quad 作为\mbox{练习~26.2.H} 实战的一个例子，我们考虑交于一点的两个平面这一例子，沿用\mbox{练习~26.1.G} 的记号。\footnote{原书此处作「two planes meeting a point」，脱「at」，当为「meeting at a point」之误，径补。——译者注}考虑由 $(w-y,\,x-z)$ 给出的\trMorphism{} $\pi\colon X\to Y:=\mathbb{A}^2$。注意 $Y$ 是正则的（\S13.3.6）。但 $X$ 不是\trCM{}的（\mbox{练习~26.1.G}），且 $\pi$ 不是平坦的（\mbox{练习~24.3.I}），于是由\mbox{练习~26.2.H}，你可以用其中任意一条去证明另一条。

\textbf{26.2.13. \trCM{}性的花哨性质。}

我们再列举\trCM{}性的几条附加性质，不加证明。它们值得一看，但我们不会用到。

设 $A$ 是诺特环，且是某个\trCM{}环的商环，则 $\Spec A$ 中\trCM{}点构成的\trLocus{}是开的；见 [Mat2, Ex. 24.2]。与此相关：

\textbf{26.2.14. 事实（参看\mbox{事实~13.8.5}）。}\ ——\ \emph{\trCM{}\trNoether{}局部环在任一\trPrimeIdeal{}处的局部化仍是\trCM{}的。}

（直接证明见 [Stacks, tag 00NB]、[E, Prop. 18.8]、[Mat2, Thm. 17.3(iii)] 或 [Mat1, Thm. 30]。）特别地，对诸簇而言，\trCM{}性可以在闭点处检验（\mbox{练习~5.3.F}）。利用「\trCM{}性在局部化下保持」这一事实，还可以很快证明：\trCM{}局部环都是\trCatenary{}的（见 [E, Cor. 18.10]、[Mat1, Thm. 31(ii)] 或 [Mat2, Thm. 17.9]）。环 $A$ 是\trCM{}的，当且仅当 $A[x]$ 是\trCM{}的（[E, Prop 18.9]、[Mat2, Thm 17.7]、[Mat1, Thm. 33]），当且仅当 $A[[x]]$ 是\trCM{}的 [Mat2, p. 137]。局部环是\trCM{}的，当且仅当它的\trCompletion{}是\trCM{}的（[E, Prop. 18.8]、[Mat2, Thm 17.5]）。

至于这些以及许多其他令人愉快的事实，还有一般地对\trCM{}性的一份透彻入门，见 [BruH]。


\section{\trSerre{}的 $R_1+S_2$ 正规性判别法}

\trDepth{}的观念给出了一个有用的\trNormality{}判别法，它归功于\trSerre{}。

\textbf{26.3.1.}\ \emph{\trhDefinition{}。}\quad 设 $A$ 是一个\trNoether{}环，而 $k\in\mathbb{Z}^{\ge0}$。我们说 $A$ 具有\textbf{性质 $R_k$}（即 $A$ \textbf{在\trCodimension{} $\le k$ 处正则}（regular in codimension $\le k$），或者更不严格地说，$A$ \textbf{在\trCodimension{} $k$ 处正则}），若对每个\trCodimension{}至多为 $k$ 的\trPrimeIdeal{} $\mathfrak{p}\subset A$，$A_{\mathfrak{p}}$ 都是正则的。（在\trCodimension{} $1$ 处正则的重要性，我们已经屡次见到。）于是，一个\trNoether{}环是正则的，当且仅当它对一切 $k$ 都具有性质 $R_k$。\footnote{原书此处作「regular and only if」，当为「regular if and only if」之误（漏排 if），径改。——译者注}

我们说 $A$ 具有\textbf{性质 $S_k$}，若对每个\trPrimeIdeal{} $\mathfrak{p}\subset A$，\trLocalRing{} $A_{\mathfrak{p}}$ 的\trDepth{}都至少为 $\min(k,\dim A_{\mathfrak{p}})$——「诸局部环直到\trCodimension{} $k$ 为止都是 \trCM{}的，而在那之后其\trDepth{}都至少为 $k$」。于是，一个\trNoether{}环是 \trCM{}的，当且仅当它对一切 $k$ 都具有性质 $S_k$。

\textbf{26.3.A.（容易练习）}\quad 注意：一个\trNoether{}环 $A$ 总是平凡地具有性质 $S_0$。证明：一个\trNoether{}环 $A$ 具有性质 $R_0$，若 $\Spec A$ 是「一般既约的」（generically reduced）：即在其每个\trIrreducibleComponent{}的\trGenericPoint{}处既约。\footnote{原书此处作「has property \(R_0\) if Spec A is generically reduced」，据\mbox{练习~26.3.C} 疑当为「if and only if」，兹照录原书。——译者注}证明：一个\trNoether{}环 $A$ 具有性质 $S_1$，当且仅当 $\Spec A$ 没有\trEmbeddedPoint{}。（可能的提示：\mbox{练习~26.1.A}。）

\textbf{26.3.B.（练习：$S_2$ 的重要初等解读）}\quad 设 $A$ 是一个\trNoether{}\trIntegralDomain。证明：$A$ 具有性质 $S_2$，当且仅当对一切非零的 $g\in A$，$A/(g)$ 都没有\trEmbeddedPrime{}。

\textbf{26.3.C.（练习）}\quad 证明：一个\trNoether{}环 $A$ 没有幂零元（即 $\Spec A$ 既约），当且仅当它同时具有性质 $R_0$ 与 $S_1$。（提示：证明一个\trNoether{}概形是既约的，当且仅当它的每个\trIrreducibleComponent{}都一般既约，并且它没有\trEmbeddedPoint{}。）

把\mbox{练习~26.3.C} 中的下标各加一，便得到\trSerre{}的判别法。

\textbf{26.3.2. \trhTheorem{}（\trSerre{}正规性判别法）。}\ ——\ \emph{一个\trNoether{}环 $A$ 是\trNormal{}的，当且仅当它具有性质 $R_1$ 与 $S_2$。}

（回想 \S5.4.1：一个\trNoether{}环 $A$ 称为\trNormal{}的，若它在一切\trPrimeIdeal{} $\mathfrak{p}$ 处的\trLocalization{} $A_{\mathfrak{p}}$ 都是\trIntegrallyClosed\trIntegralDomain。）于是，\trNormality{}的失效可能有两个原因：它可能是 $R_1$ 的失效（这是我们已然知道的，来自\mbox{定理~13.5.8} 中 (a) 与 (g) 的等价性），也可能是 $S_2$ 更为微妙的失效。满足 $R_1$ 但不满足 $S_2$ 的簇的例子，见\mbox{例~26.3.4}（交于一点的两个平面）与\mbox{练习~26.3.F}（\trPinchedPlane）。

\textbf{26.3.D.（容易练习）}\quad 证明\trSerre{}正规性判别法（\mbox{定理~26.3.2}）。提示：仔细考察\mbox{命题~13.5.24} 与\mbox{命题~13.5.26}。

\textbf{26.3.3. 应用。}\quad 首先，\trSerre{}正规性判别法蕴含：\trCM{}概形是\trNormal{}的，当且仅当它们在\trCodimension{} $1$ 处正则。于是，要检验 $\mathbb{P}^n_k$ 中（或者更一般地，在任何光滑簇中）的超曲面（或更一般地，\trRegularEmbedding）的\trNormality，只须检验它们的\trSingularLocus{}的\trCodimension{}大于 $1$。（这些陈述之所以为真，你应当把其中的细节想清楚。）特别地，这给出了\mbox{练习~5.4.I}(b) 与 (c) 的一个新（更为复杂的）证明。

\textbf{26.3.E.（练习：概念演练）}\quad 证明：二维\trNormal\trVariety{}都是 \trCM{}的。

\textbf{26.3.4.}\ \emph{例：交于一点的两个平面。}\quad \mbox{练习~26.1.G} 中的簇 $X$（交于一点的两个平面）不是\trNormal{}的（为什么？），但它在原点之外处处正则。这就蕴含：$X$ 不具有性质 $S_2$（从而不是 \trCM{}的），而无须\mbox{练习~26.1.G} 的那些代数操作。

我们已然知道这个例子不是 \trCM{}的，但同样的想法可以表明：\trPinchedPlane{} $\Spec k[x^3,x^2,xy,y]$（出现在\mbox{练习~13.5.H} 中）不是 \trCM{}的。由于该环的「外在」描述，很难用别的方式来做到这一点。

\textbf{26.3.F.（练习：\trPinchedPlane{}不是 \trCM{}的）}\quad 设 $A$ 是 $k[x,y]$ 的子环 $k[x^3,x^2,xy,y]$。证明：$A$ 不是 \trCM{}的。提示：\mbox{练习~13.5.H} 已表明 $A$ 不是整闭的。

\textbf{26.3.5.}\ \emph{\trRegular\trLocalRing{}是整闭的。}\quad \trSerre{}正规性判别法可以用来证明：正则局部环是整闭的——而无须绕道困难的\mbox{事实~13.8.5}（\mbox{Auslander--Buchsbaum 定理}），即它们都是\trUFD。正则局部环都是 \trCM{}的（\S26.2.5），并且（由\mbox{练习~26.3.C}）因它们是整环（\mbox{定理~13.2.11}）而在\trCodimension{} $0$ 处正则，于是我们只须证明：正则局部环在\trCodimension{} $1$ 处正则。我们可以援引另一条困难的\mbox{事实~13.8.2}，即正则局部环的\trLocalization{}仍是正则的；但至少对于\trPerfectField{}上有限生成代数的\trLocalization，我们已经证明了这一点（见\mbox{定理~13.8.3} 与\mbox{练习~21.6.B}）。

\textbf{26.3.6.}\ \emph{\trhCaveat{}。}\quad 正如下面的练习所清楚表明的，条件 $S_2$ 是对\emph{一切}素理想施加的条件，而不仅仅是\trCodimension{}至多为 $2$ 的那些素理想。

\textbf{26.3.G.（练习）}\quad 举出一个满足 $R_1$、并且在一切\trCodimension{}至多为 $2$ 的点处都是 \trCM{}的簇的例子，它不是\trNormal{}的。

\textbf{26.3.7.}\ \extrahard\ \emph{射影正规性。}\quad 我们不会使用射影正规性这一概念，但仍要提它一笔，因为许多读者可能会与它打交道，而且把它介绍出来很容易。它是闭嵌入 $\pi\colon X\hookrightarrow\mathbb{P}^n_k$ 的一条性质。令 $\mathscr{L}=\pi^*\OO(1)$ 是给出这一闭嵌入的\trLineBundle（\S15.2），并且为方便起见，令 $R(\mathscr{L})_\bullet:=\bigoplus_{m\ge0}\Gamma(X,\mathscr{L}^{\otimes m})$ 为「$\mathscr{L}$ 的一切幂的截面」的\trGradedRing。于是我们说（定义 1）：$X\hookrightarrow\mathbb{P}^n_k$ 是\textbf{射影正规的}（projectively normal），若 $X$ 是正规的，并且对一切 $m\ge0$，$\mathrm{H}^0(\mathbb{P}^n_k,\OO(m))\to\mathrm{H}^0(X,\OO_X(m))$ 都是满射。换言之，对一切 $m\ge0$，$\mathscr{L}^{\otimes m}$ 的一切截面都「来自」多项式；不存在不「源出」于 $\mathrm{H}^0(X,\mathscr{L})$ 的「新」截面。等价地（定义 2）：经由 $\mathscr{L}=\pi^*\OO(1)$，$R(\mathscr{L})_\bullet$ 是多项式环 $k[x_0,\dots,x_n]$ 的一个商。你很快就能证明：这等价于如下陈述（定义 3）：（$X$ 是正规的，并且）对一切 $m\ge0$ 都有 $\mathrm{H}^1(\mathbb{P}^n_k,\mathscr{I}_X(m))=0$；也等价于如下陈述（定义 4）：$X$ 是正规的，$R(\mathscr{L})_\bullet$ 由次数 $1$ 生成，而且 $X\hookrightarrow\mathbb{P}^n_k$ 是由某个完全线性系给出的嵌入（假定它是非\trDegenerate{}的）。

\textbf{26.3.H.（练习）}\quad 证明：$\mathbb{P}^n_k$ 中的\trCompleteIntersection{}是射影正规的，当且仅当它们是正规的，当且仅当它们在\trCodimension{} $1$ 处正则。

证明 $X\hookrightarrow\mathbb{P}^n_k$ 是正规的，当且仅当其\trAffine\trCone{}是正规的，这并不太难；见 [Ha1, Ex.~II.5.14]。

